\documentclass[../../../include/open-logic-section]{subfiles}

\begin{document}
\olfileid{sth}{ordinals}{zfm}
\olsection{$\ZFminus$: معلم مهم}

أما تسويغ الاستبدال، وهل يمكن تسويغه أصلًا، فمسألة ليست باليسيرة؛
نؤخر بحثها إلى~\olref[replacement][]{chap}. لكن إضافته أوصلتنا إلى
معلم آخر مهم، إذ استوفت لنا مسلّمات النظرية~$\ZFminus$. وتفصيلها الآتي.

\begin{defn}
تتألف مسلّمات~$\ZFminus$ من الامتدادية، والاتحاد، والأزواج، ومجموعات
القوى، واللانهاية، ومن جميع حالات مخططي الفصل والاستبدال.
فـ$\ZFminus$، بعبارة أخرى، هي~$\Zminus$ مضافًا إليها الاستبدال.
\end{defn}
\noindent
والاسم اختصار لنظرية \emph{زرملو--فرينكل} للمجموعات، \emph{منقوصًا}
منها شيء يأتي ذكره. وإنما ذُكر فرينكل معه لنسبة صوغ الاستبدال إليه
في~\citeyear{Fraenkel1922}؛ على أن أول صياغة دقيقة وضعها
\citet{Skolem1922}.

\end{document}
